\begin{abstract}
We prove that every finite bimolecular weakly reversible stochastic
mass-action network with one linkage class is positive recurrent on each
closed communicating class, for every positive choice of rate constants.
This removes the hypothesis in the 2020 theorem of Anderson, Cappelletti,
and Kim that each species occurs in a unary or double complex.  The proof
marks the target complex of the actual reaction channel that most recently
fired and applies a log-factorial potential to the population remaining after
that target is removed.  Its increment is exactly a target/source
falling-factorial ratio.  Finite target-following paths propagate the negative
reward of a rare terminal source, while a normalized-log compactification and
an exhaustive bimolecular top-complex alternative make that terminal source
available along every divergent sequence unless a linear stoichiometric
invariant already bounds the sequence.
\end{abstract}

\noindent\textbf{Keywords.}
Stochastic reaction network; chemical reaction network; positive recurrence;
continuous-time Markov chain; mass-action kinetics; weak reversibility;
Foster--Lyapunov method.

\medskip
\noindent\textbf{2020 Mathematics Subject Classification.}
Primary 60J27; secondary 60J74, 92C42.

\section{Introduction}\label{sec:intro}

Stochastic reaction networks are continuous-time Markov-chain models for
interacting populations whose transition rates are determined by a finite
reaction graph.  They arise in chemical kinetics, systems biology, ecology,
and related population models; see, for example,
\citet{AndersonKurtz2015}.  A basic stability question is whether each closed
irreducible component has a stationary probability distribution, or
equivalently whether its states have finite mean return times.

\citet{AndersonKim2018} formulated the positive-recurrence conjecture that
weak reversibility should force this stability for stochastic mass-action
systems.  They established several structural sufficient conditions.
\citet{AndersonCappellettiKim2020} subsequently proved the binary,
one-linkage case under the additional condition
\begin{equation}\label{eq:ACK-condition}
  \{S_i,2S_i\}\cap\mathcal C\ne\varnothing
  \qquad\text{for every species }S_i.
\end{equation}
Thus each species had to occur in a pure unary or pure double complex.  Other
progress includes complex-balanced product-form systems
\citep{AndersonCraciunKurtz2010}, strongly endotactic structures
\citep{AndersonCappellettiKimNguyen2020}, one-dimensional classifications
\citep{WiufXu2023}, and nonexplosion of all second-order weakly reversible
systems \citep{Xu2026}.  The May 2026 revision of \citet{Xu2026} still records
the general bimolecular positive-recurrence problem as open.

The present paper removes \eqref{eq:ACK-condition} while retaining the binary
and one-linkage assumptions.  The exact theorem is stated in
Theorem~\ref{thm:main}.  It does not address multiple linkage classes or
complexes of molecularity above two, and therefore does not settle the full
positive-recurrence conjecture.

A small motivating example is the directed cycle
\begin{equation}\label{eq:stress}
  0\longrightarrow A+B\longrightarrow B\longrightarrow0.
\end{equation}
Species \(A\) occurs in neither \(A\) nor \(2A\), so
\eqref{eq:ACK-condition} fails.  At the boundary state \((n,0)\), only the
reaction \(0\to A+B\) is enabled and the one-step drift of total population
is positive.  The restoring effect is delayed: the created \(B\) enables a
fast reaction consuming \(A\).  This is representative of why a one-step
radial Lyapunov function is not the right universal object.

The proof instead uses four ideas.  First, we augment the embedded jump chain
by marking the actual reaction channel that fired and retaining its target
complex.  Second, after removing that target from the population, the
log-factorial potential has the exact increment
\[
  \log\frac{(x)_t}{(x)_s}
\]
when the carried target is \(t\) and the next source is \(s\).  Following the
carried target therefore costs exactly zero.  Third, a finite episode follows
a fixed directed path from the carried target toward a selected terminal
complex.  A sharp scalar envelope shows that a terminal source probability
tending to zero forces the complete episode drift to tend to \(-\infty\), no
matter how the intermediate source rates are separated.  Fourth, a
normalized-log compactification of an arbitrary divergent sequence gives an
exhaustive top-complex case split.  Bimolecularity implies either the
existence of a higher-weight source enabled over a lower terminal complex or
a reaction-wise linear invariant that contradicts divergence inside a fixed
communicating class.  A finite-family random-time drift argument and a finite
trace-chain construction complete the proof.

\section{Reaction networks and the stochastic model}\label{sec:model}

Let \(\mathcal S=\{S_1,\dots,S_d\}\) be a finite species set.  A
\emph{complex} is a vector \(y=(y_1,\dots,y_d)\in\mathbb N_0^d\), identified
with the formal sum \(\sum_i y_iS_i\).  Its molecularity is
\(|y|=\sum_i y_i\).  A \emph{reaction channel} is a labelled ordered pair
\(r:y_r\to y_r'\), together with a positive rate constant \(\kappa_r\).
Parallel channels are allowed.  The reaction graph has the complexes as
vertices and the reaction channels as directed edges.  Its linkage classes
are the connected components of the underlying undirected graph.  A network
is \emph{weakly reversible} when each linkage class is strongly connected.
With one linkage class, weak reversibility is simply strong connectivity of
the entire complex graph.

We call the network \emph{binary} or \emph{bimolecular} when
\begin{equation}\label{eq:binary}
  |y|\le2\qquad\text{for every }y\in\mathcal C.
\end{equation}
Because every complex in a weakly reversible graph is a reaction source, this
agrees here with the usual second-order/bimolecular terminology.

For \(x\in\mathbb N_0^d\), define the multivariate falling factorial
\begin{equation}\label{eq:falling}
  (x)_y=
  \begin{cases}
  \displaystyle\prod_{i=1}^d\frac{x_i!}{(x_i-y_i)!},&x\ge y,\\[1ex]
  0,&\text{otherwise}.
  \end{cases}
\end{equation}
The stochastic mass-action propensity of channel \(r\) is
\(\lambda_r(x)=\kappa_r(x)_{y_r}\), and the formal generator is
\begin{equation}\label{eq:generator}
  \mathcal Lf(x)=\sum_r\kappa_r(x)_{y_r}
  \bigl[f(x+y_r'-y_r)-f(x)\bigr].
\end{equation}
The convention using products of binomial coefficients is equivalent after
rescaling each positive rate constant by a fixed factorial factor.

A subset \(\Gamma\subseteq\mathbb N_0^d\) is a \emph{closed communicating
class} when its states communicate through enabled channels and no enabled
channel exits \(\Gamma\).  The chain restricted to \(\Gamma\) is
irreducible.  It is \emph{positive recurrent} when one, equivalently every,
state has finite expected positive return time.  For a nonexplosive
irreducible countable-state CTMC, this is equivalent to the existence of a
stationary probability distribution.

A channel \(y\to y\) has zero population increment and contributes zero to
\eqref{eq:generator}; discard it.  Exact parallel channels with the same
source and target may be combined by adding their rates.  Channels having the
same population displacement but different sources or targets remain
distinct, because the proof marks the actual channel.

Every finite closed communicating class is positive recurrent: its finite
irreducible jump chain has a stationary probability distribution, all state
rates are finite, and the associated finite-state CTMC is nonexplosive.  If
the complex graph has one vertex, all closed classes are absorbing
singletons.  If all complexes have the same molecularity, every reaction
preserves \(|x|_1\), so each communicating class lies in a finite shell.  We
therefore fix henceforth an infinite closed irreducible class \(\Gamma\), a
strongly connected graph with at least two complexes, and no self channels.
Every complex then has a genuine outgoing channel.  Write
\begin{equation}\label{eq:kbar}
  \bar\kappa_y=\sum_{r:y_r=y}\kappa_r,
  \qquad
  \bar\kappa_- =\min_{y\in\mathcal C}\bar\kappa_y>0,
  \qquad
  \bar\kappa_+ =\max_{y\in\mathcal C}\bar\kappa_y.
\end{equation}

\begin{theorem}[Binary one-linkage positive recurrence]\label{thm:main}
Let a finite stochastic mass-action reaction network be weakly reversible,
have one linkage class, and satisfy \eqref{eq:binary}.  For every positive
rate vector, every closed communicating class of its population CTMC is
positive recurrent.
\end{theorem}

\begin{corollary}[Stationary probability]\label{cor:stationary}
Under the hypotheses of Theorem~\ref{thm:main}, every closed communicating
class \(\Gamma\) admits a stationary probability distribution
\(\pi_\Gamma\).  It is unique on \(\Gamma\).
\end{corollary}

\section{The marked target-augmented chain}\label{sec:marked}

We first define the embedded chain without assuming nonexplosion.  At a
nonabsorbing population state \(x\), sample the actual labelled reaction
channel \(r\) with probability
\begin{equation}\label{eq:channel-kernel}
  \frac{\kappa_r(x)_{y_r}}{\Lambda(x)},
  \qquad
  \Lambda(x)=\sum_q\kappa_q(x)_{y_q}.
\end{equation}
This discrete chain can be constructed independently of its holding times.
After the sampled channel \(r\) fires, record its target \(t=y_r'\).  The mark
is therefore obtained from the actual channel and is never inferred from the
population displacement.  Once exact parallel channels with the same source
and target have been aggregated, retaining the target is the only part of the
previous channel needed below.  A reachable augmented state is denoted
\((x,t)\), where \(x\) is the post-jump population.  If the next sampled
channel is \(q:s\to u\), then
\begin{equation}\label{eq:aug-transition}
  (x,t)\longmapsto(x-s+u,u).
\end{equation}
The channel law depends on \(x\), not on the historical mark, so the augmented
process is Markov, and projection onto \(x\) is the ordinary population
embedded chain.

Let \(\widetilde\Gamma\) be the augmented states reachable after genuine
jumps within \(\Gamma\).  It is closed under \eqref{eq:aug-transition}.  It is
also irreducible.  Indeed, given \((x,t),(x',t')\in\widetilde\Gamma\), choose
a channel \(s'\to t'\) and predecessor population
\(z=x'-t'+s'\in\Gamma\) witnessing reachability of the second pair.
Population irreducibility supplies a positive-probability marked-channel path
from \(x\) to \(z\); appending \(s'\to t'\) reaches \((x',t')\).

The carried target is always enabled.  If the preceding channel was
\(s\to t\), then its post-jump state has the form \(x=z-s+t\ge t\).  Put
\begin{equation}\label{eq:residual}
  r=x-t\in\mathbb N_0^d,
  \qquad
  V(x,t)=\sum_{i=1}^d\log((x_i-t_i)!)=\sum_i\log(r_i!).
\end{equation}
We use \(\log(0!)=0\).  Thus \(V\ge0\).  Every sublevel set of \(V\) in
\(\widetilde\Gamma\) is finite, because \(\log(n!)\to\infty\), each residual
coordinate is bounded on a sublevel, and the target set is finite.

\begin{lemma}[Exact target/source identity]\label{lem:identity}
At \((x,t)\), if the next marked channel has source \(s\) and target \(u\),
then
\begin{equation}\label{eq:identity}
  V(x-s+u,u)-V(x,t)=\log\frac{(x)_t}{(x)_s}.
\end{equation}
In particular, a channel sourced at the carried target has zero increment.
\end{lemma}

\begin{proof}
The new residual is \((x-s+u)-u=x-s\).  Since both \(s\) and \(t\) are
enabled,
\[
 \exp(\Delta V)
 =\prod_i\frac{(x_i-s_i)!}{(x_i-t_i)!}
 =\frac{(x)_t}{(x)_s}.
\]
\end{proof}

This identity gives a precise ``complete-credit'' interpretation: a product
complex remains credited in full until the chain chooses a different source.
The interpretation is not used in place of the algebraic identity.

For enabled source complexes define
\begin{equation}\label{eq:source-prob}
  p_x(y)=\frac{\bar\kappa_y(x)_y}{\Lambda(x)}.
\end{equation}
Set \(p_x(y)=0\) for disabled sources, and use \(0\log0=0\) in entropy
expressions.  Unless explicitly stated otherwise, sums in this subsection
range only over enabled source complexes.  Because the carried target is
enabled, \(p_x(t)>0\).  The expected one-jump increment is
\begin{equation}\label{eq:d-def}
  d(x,t)=\sum_{s:\,x\ge s}p_x(s)
  \log\frac{(x)_t}{(x)_s}.
\end{equation}

\begin{lemma}[Source-probability bound]\label{lem:entropy}
For every reachable augmented state,
\begin{equation}\label{eq:entropy-identity}
\begin{split}
  d(x,t)
  ={}&\log p_x(t)-\sum_{s:\,x\ge s} p_x(s)\log p_x(s)\\
    &+\sum_{s:\,x\ge s} p_x(s)\log\bar\kappa_s
      -\log\bar\kappa_t,
\end{split}
\end{equation}
and consequently
\begin{equation}\label{eq:entropy-bound}
  d(x,t)\le\log p_x(t)+C_0,
  \qquad
  C_0=\log|\mathcal C|+
  \log\frac{\bar\kappa_+}{\bar\kappa_-}.
\end{equation}
\end{lemma}

\begin{proof}
For each enabled \(s\),
\((x)_s=p_x(s)\Lambda(x)/\bar\kappa_s\).  Substitution into
\eqref{eq:d-def} gives \eqref{eq:entropy-identity}.  The entropy of a
probability vector supported on at most \(|\mathcal C|\) points is at most
\(\log|\mathcal C|\), and the last line contributes at most
\(\log(\bar\kappa_+/\bar\kappa_-)\).  The zero complex creates no exception,
since \((x)_0=1\).
\end{proof}

\section{Propagation along target-following paths}\label{sec:paths}

For every ordered pair \((t,c)\in\mathcal C^2\), fix a simple directed path
of marked channels
\begin{equation}\label{eq:path}
 t=y_0\longrightarrow y_1\longrightarrow\cdots
 \longrightarrow y_L=c,
 \qquad 0\le L\le|\mathcal C|-1.
\end{equation}
Strong connectivity supplies the path.  Starting from \((r+t,t)\), define
the \((t,c)\)-episode as follows.  At phase \(y_k\), \(k<L\), continue only
when the exact designated channel \(y_k\to y_{k+1}\) fires; after any other
channel, stop immediately.  If the path reaches \(c\), take one final
ordinary jump and stop.  When \(c=t\), the path has length zero and the
episode is just that final jump.

This is a stopping time for the marked embedded filtration.  On designated
edges the lifted populations are
\begin{equation}\label{eq:lift}
  r+y_0\longrightarrow r+y_1\longrightarrow\cdots\longrightarrow r+c.
\end{equation}
Every designated source is literally present, the residual remains \(r\),
and closedness of \(\Gamma\) keeps the entire lift in the class.  The episode
has between one and \(|\mathcal C|\) jumps and only finitely many terminal
branches.

For the designated channel at phase \(k\), put
\begin{equation}\label{eq:qk}
 q_k=\frac{\kappa_{y_k\to y_{k+1}}}{\bar\kappa_{y_k}}>0,
 \qquad p_k=p_{r+y_k}(y_k).
\end{equation}
Let \(J_k(r)\) be the expected remaining \(V\)-increment from phase \(k\).

\begin{lemma}[Exact episode recursion]\label{lem:recursion}
The episode rewards satisfy
\begin{equation}\label{eq:recursion}
  J_L(r)=d(r+c,c),
  \qquad
  J_k(r)=d(r+y_k,y_k)+q_kp_kJ_{k+1}(r).
\end{equation}
\end{lemma}

\begin{proof}
At phase \(k\), the first jump has expected increment
\(d(r+y_k,y_k)\).  Continuation occurs exactly when the designated channel
fires, with probability \(q_kp_k\).  That channel is sourced at the carried
target and has zero immediate increment by Lemma~\ref{lem:identity}; the
remaining conditional expectation is \(J_{k+1}(r)\).
\end{proof}

The following scalar calculation controls all possible intermediate source
probabilities.

\begin{lemma}[Scalar envelope]\label{lem:envelope}
For \(q>0\) and \(M\in\mathbb R\), define
\[
  F_q(M)=\sup_{0<p\le1}\{\log p+C_0+qpM\}.
\]
Then
\begin{equation}\label{eq:F-piecewise}
  F_q(M)=
  \begin{cases}
    C_0+qM,&M\ge-1/q,\\
    C_0-1-\log(-qM),&M<-1/q.
  \end{cases}
\end{equation}
Consequently, if \(M_n\to-\infty\), then \(F_q(M_n)\to-\infty\).
\end{lemma}

\begin{proof}
The objective has derivative \(p^{-1}+qM\) and second derivative
\(-p^{-2}<0\).  If \(M\ge-1/q\), its derivative at \(p=1\) is nonnegative,
so concavity places the maximum at \(p=1\).  If \(M<-1/q\), the unique
maximizer is \(p=(-qM)^{-1}\in(0,1)\), which gives the second line of
\eqref{eq:F-piecewise}.  The final assertion follows immediately.
\end{proof}

\begin{proposition}[Propagation of terminal drift]\label{prop:propagation}
Fix a target-following path \eqref{eq:path}.  If along a sequence of residuals
\(r^{(n)}\) one has
\[
  p_{r^{(n)}+c}(c)\longrightarrow0,
\]
then the expected increment of the complete \((t,c)\)-episode tends to
\(-\infty\).
\end{proposition}

\begin{proof}
At the terminal phase, Lemma~\ref{lem:entropy} gives
\(J_L\le\log p_{r+c}(c)+C_0\to-\infty\).  If \(J_{k+1}\le M\),
Lemmas~\ref{lem:entropy} and \ref{lem:recursion} give
\[
  J_k\le\sup_{0<p\le1}\{\log p+C_0+q_kpM\}=F_{q_k}(M).
\]
Apply Lemma~\ref{lem:envelope} backward through the finite path.
\end{proof}

The proof preserves every reaction-rate monomial: no edge is replaced and no
unrelated rates are compared.  Arbitrarily many nested source-rate scales are
absorbed by the finite scalar recursion.

\section{Logarithmic compactification and the top-complex alternative}
\label{sec:compact}

We now show that every divergent augmented sequence supplies the terminal
rarity required by Proposition~\ref{prop:propagation}, unless a linear
stoichiometric invariant already makes such divergence impossible.

Take a sequence \((x^{(n)},t^{(n)})\in\widetilde\Gamma\) leaving every finite
set.  Pass to a subsequence with \(t^{(n)}=t\), and set
\(r^{(n)}=x^{(n)}-t\).  A diagonal extraction makes every coordinate of
\(r^{(n)}\) either one fixed nonnegative integer or divergent to infinity.
Let \(I\) be the set of divergent coordinates and put
\begin{equation}\label{eq:Rn}
  R_n=\sum_{i\in I}\log(r_i^{(n)}+1).
\end{equation}
After a further subsequence, the limits
\begin{equation}\label{eq:weights}
  w_i=\lim_{n\to\infty}
  \frac{\log(r_i^{(n)}+1)}{R_n}
  \quad(i\in I),
  \qquad w_i=0\quad(i\notin I)
\end{equation}
exist.  Thus \(w\ge0\) and \(\sum_iw_i=1\).  A divergent coordinate may
have \(w_i=0\); it remains in \(I\), so all slower tiers are retained.

\begin{lemma}[Falling-factorial asymptotics]\label{lem:ff-asymptotic}
Let \(c\) and \(y\) be fixed complexes, and suppose \(y\) is enabled at
\(r^{(n)}+c\) for all sufficiently large \(n\).  Then
\begin{equation}\label{eq:ff-asymptotic}
  \log(r^{(n)}+c)_y=R_n\,w\cdot y+o(R_n).
\end{equation}
\end{lemma}

\begin{proof}
For \(y_i=1\) and a divergent coordinate,
\(\log(r_i^{(n)}+c_i)=\log(r_i^{(n)}+1)+o(R_n)\); a fixed coordinate
contributes \(O(1)\).  For \(y_i=2\) and \(r_i^{(n)}\to\infty\),
\[
 \log\bigl((r_i^{(n)}+c_i)(r_i^{(n)}+c_i-1)\bigr)
 =2\log(r_i^{(n)}+1)+o(R_n).
\]
If that coordinate is fixed, enabledness makes its contribution finite and
hence \(O(1)\).  Summing the coordinate contributions proves the claim.
\end{proof}

For a complex \(y\), write \(h(y)=w\cdot y\), and define
\begin{equation}\label{eq:top}
  a=\max_{y\in\mathcal C}h(y),
  \qquad
  T=\{y\in\mathcal C:h(y)=a\}.
\end{equation}
The next lemma is the load-bearing use of bimolecularity.

\begin{lemma}[Top-complex availability or invariant]\label{lem:top}
Along the divergent sequence, at least one of the following conclusions is
available:

\begin{enumerate}[label=\textup{(\roman*)},leftmargin=2.5em]
\item a reaction-wise linear stoichiometric invariant has value tending to
      infinity along the sequence, contradicting membership in the fixed
      communicating class; or
\item there are fixed complexes \(s,c\in\mathcal C\) with
      \(h(s)>h(c)\) such that both \(s\) and \(c\) are enabled at
      \(r^{(n)}+c\) for all sufficiently large \(n\).
\end{enumerate}
\end{lemma}

\begin{proof}
The following cases are exhaustive.

\begin{center}
\begin{tabular}{p{0.31\linewidth}p{0.59\linewidth}}
\toprule
Configuration & Conclusion \\
\midrule
Every complex is top & The nonnegative invariant \(w\cdot X\). \\
A top complex contains two divergent particles & That top complex is enabled over every lower terminal. \\
Every top complex contains one divergent particle & Reduce to the four subcases below using the set \(K\) of top species. \\
\bottomrule
\end{tabular}
\end{center}

If every complex is top, then \(w\cdot(y'-y)=0\) for every channel.
Therefore \(w\cdot X\) is a nonnegative linear invariant.  A coordinate with
positive weight diverges, so the invariant value tends to infinity, which is
impossible inside one communicating class.

Assume now that \(T\ne\mathcal C\).  Then \(a>0\).  If a top complex \(s\)
contains two particles whose coordinates belong to \(I\), bimolecularity
forces \(s\) to contain no bounded-coordinate particle.  Both required
particles are eventually present in \(r^{(n)}\), including the case
\(s=2S_i\).  Hence \(s\) is enabled at \(r^{(n)}+c\) for every lower
\(c\notin T\), and \(h(s)>h(c)\).

It remains to suppose that every top complex contains exactly one divergent
particle.  Let \(K\subseteq I\) be the species appearing in top complexes.
If \(i\in K\), the remaining possible particle in that top complex has
weight zero, and therefore \(w_i=a\).  No complex contains two
\(K\)-particles, since its weight would be at least \(2a>a\).  Moreover,
\begin{equation}\label{eq:top-iff}
  y\in T\quad\Longleftrightarrow\quad
  q_K(y):=\sum_{i\in K}y_i=1.
\end{equation}
The forward implication follows from the definition of \(K\).  Conversely,
one \(K\)-particle already contributes \(a\); maximality forces all other
weight to be zero and makes the complex top.

There are four subcases.

\begin{enumerate}[label=\textup{3\alph*.},leftmargin=3em]
\item If every complex has \(q_K=1\), then
      \(M_K=\sum_{i\in K}X_i\) is a nonnegative invariant.  It diverges along
      the sequence, a contradiction.
\item If a unary top complex \(S_i\), \(i\in K\), exists, it is enabled over
      every lower terminal complex, giving conclusion (ii).
\item Otherwise every top complex has the form \(S_i+D\), where
      \(D\notin I\); call \(D\) a service species.  If a lower complex
      \(c\) contains such a \(D\), the corresponding top source \(S_i+D\)
      is enabled at \(r^{(n)}+c\), giving conclusion (ii).
\item If no lower complex contains any service species, let \(\mathcal D\)
      be the set of distinct service species.  The signed linear functional
      \begin{equation}\label{eq:signed-invariant}
        M_K-\sum_{D\in\mathcal D}X_D
      \end{equation}
      has value zero on every complex: each top complex contains one
      \(K\)-particle and one service particle, whereas a lower complex
      contains neither.  It is therefore reaction-wise invariant.  Its
      positive part diverges, while every service coordinate lies outside
      \(I\) and is fixed along the subsequence.  Its value tends to infinity,
      contradicting the fixed communicating class.
\end{enumerate}

This exhausts all binary complexes.  We reserve ``conservation law'' for the
nonnegative invariants above; \eqref{eq:signed-invariant} is a signed linear
stoichiometric invariant.
\end{proof}

The invariant contradiction is exact: if \(\ell\cdot(y'-y)=0\) on every
channel, then \(\ell\cdot x\) is constant along every path in \(\Gamma\).
In the signed branch, the negative-coordinate populations are fixed along
the extracted sequence, so the invariant really does tend to \(+\infty\).
A species absent from every complex is itself reaction-wise invariant; thus an
absent coordinate cannot diverge along a fixed communicating class.

\begin{lemma}[Vanishing terminal source probability]\label{lem:terminal}
Under conclusion \textup{(ii)} of Lemma~\ref{lem:top},
\begin{equation}\label{eq:terminal-zero}
  p_{r^{(n)}+c}(c)\longrightarrow0.
\end{equation}
\end{lemma}

\begin{proof}
Both \(s\) and \(c\) are enabled.  Lemma~\ref{lem:ff-asymptotic} and the
strict gap \(h(s)>h(c)\) give
\[
  \frac{(r^{(n)}+c)_s}{(r^{(n)}+c)_c}\longrightarrow\infty.
\]
Consequently,
\[
 p_{r^{(n)}+c}(c)
 \le
 \frac{\bar\kappa_c(r^{(n)}+c)_c}
      {\bar\kappa_s(r^{(n)}+c)_s}
 \longrightarrow0.
\]
\end{proof}

\section{Uniform random-time drift outside a finite set}\label{sec:uniform}

For each \(c\in\mathcal C\), let \(D_c(x,t)\) denote the exact expected
increment of the fixed \((t,c)\)-episode.  Define
\begin{equation}\label{eq:K}
 K=\left\{(x,t)\in\widetilde\Gamma:
          \min_{c\in\mathcal C}D_c(x,t)>-1\right\}.
\end{equation}

\begin{proposition}[Finite and nonempty exceptional set]\label{prop:K}
The set \(K\) is finite and nonempty.
\end{proposition}

\begin{proof}
Suppose first that \(K\) were infinite.  Properness of \(V\) supplies a
sequence in \(K\) leaving every finite set.  Apply the compactification of
Section~\ref{sec:compact}.  An invariant conclusion from
Lemma~\ref{lem:top} contradicts membership in one communicating class.
Otherwise, Lemma~\ref{lem:terminal} and
Proposition~\ref{prop:propagation} supply a fixed terminal complex \(c\) for
which \(D_c(x^{(n)},t)\to-\infty\).  This contradicts the definition of
\(K\).  Thus \(K\) is finite.

To prove nonemptiness, select any augmented state \(z_0\).  The finite
nonempty sublevel set \(\{z:V(z)\le V(z_0)\}\) contains a global minimizer
\(z_*\) of \(V\).  Every endpoint of every target-following episode remains
in \(\widetilde\Gamma\), so its potential is at least \(V(z_*)\).  Hence
\(D_c(z_*)\ge0\) for every \(c\), and \(z_*\in K\).
\end{proof}

Outside \(K\), choose the lexicographically first terminal complex minimizing
\(D_c(x,t)\), before the episode begins.  This is a deterministic measurable
selector on the countable state space and is compatible with restarting
under the strong Markov property.  Its episode stopping time \(\tau\) obeys
\begin{equation}\label{eq:episode-drift}
  1\le\tau\le L:=|\mathcal C|,
  \qquad
  \mathbb E_{(x,t)}[V(Z_\tau)-V(x,t)]\le-1,
  \quad (x,t)\notin K.
\end{equation}
The coordinate overshoot in one episode is deterministically bounded by
\(2L\), because each binary reaction changes each population coordinate by
at most two.

\section{Positive recurrence}\label{sec:recurrence}

We now convert \eqref{eq:episode-drift} to a finite mean return time without
invoking an unnamed random-time Foster theorem.

Let \(Y_0=z\), and let \(Y_{n+1}\) be the endpoint of the episode selected at
\(Y_n\).  Put
\(
 \sigma_K=\inf\{n\ge0:Y_n\in K\}.
\)
For each integer \(N\), sum the conditional drift inequalities only through
the bounded index \(N\wedge\sigma_K\).  Telescoping gives
\begin{equation}\label{eq:stopped-drift}
  \mathbb E_zV(Y_{N\wedge\sigma_K})
  +\mathbb E_z(N\wedge\sigma_K)
  \le V(z).
\end{equation}
Because \(V\ge0\), monotone convergence yields
\begin{equation}\label{eq:episodes-to-K}
  \mathbb E_z\sigma_K\le V(z).
\end{equation}
The original augmented embedded chain therefore hits \(K\) in finite
expected jump count, bounded by \(L V(z)\).

The following finite trace argument turns hitting \(K\) into positive return
to one state.  Details are included because hitting a finite set is not by
itself a positive-recurrence proof.

\begin{proposition}[Finite trace-chain closure]\label{prop:trace}
The augmented embedded chain has a state with finite expected positive return
time.
\end{proposition}

\begin{proof}
Let \(T_K^+=\inf\{n\ge1:Z_n\in K\}\).  From \(k\in K\), take one ordinary
jump.  There are finitely many channel successors, and from each successor
the expected hitting time of \(K\) is finite by \eqref{eq:episodes-to-K}.
Thus \(\mathbb E_kT_K^+<\infty\) for every \(k\in K\).

Record successive visits to \(K\).  The resulting trace chain is finite and
irreducible.  Indeed, any augmented path between two points of \(K\), with
its intermediate \(K\)-visits retained, gives a positive-probability trace
path.  Fix \(k_*\in K\).  From every trace state choose a finite
positive-probability trace path to \(k_*\).  Finiteness supplies common
constants \(m<\infty\) and \(\varepsilon>0\) such that, from every trace
state, \(k_*\) is hit within the next \(m\) trace transitions with probability
at least \(\varepsilon\).  Starting from \(k_*\), first take one trace
transition and then apply this estimate in successive blocks of \(m\)
transitions.  The number of blocks before the next visit to \(k_*\) is
geometrically dominated, so the positive trace return time has finite mean.

Finally, set
\[
  B=\max_{k\in K}\mathbb E_kT_K^+<\infty.
\]
If \(M\) is the positive trace return count and \(L_j\) is the number of
original jumps in its \(j\)-th trace excursion, the strong Markov property
gives \(\mathbb E[L_j\mid\mathcal F_j]\le B\).  Tonelli's theorem therefore
yields
\[
  \mathbb E\!\left[\sum_{j=0}^{M-1}L_j\right]
  \le B\,\mathbb E M<\infty.
\]
Hence the augmented embedded chain has finite expected positive return time
to \(k_*\).
\end{proof}

Projection away from the target mark shows that the population embedded chain
has finite expected return to the population coordinate \(x_*\) of \(k_*\).
It remains to construct physical time and rule out explosion.

\begin{proposition}[Embedded-chain-to-CTMC conversion]\label{prop:ctmc}
Let a countable irreducible embedded chain be positive recurrent.  Suppose a
pure-jump construction assigns a finite rate \(\Lambda(x)>0\) at every state
and that \(\inf_x\Lambda(x)>0\).  Then the resulting minimal CTMC is
nonexplosive and positive recurrent.
\end{proposition}

\begin{proof}
Let \(N\) be a finite-mean positive return count to a recurrent embedded state
\(x_*\), and let \(H_j\) be the holding times.  Conditional on the embedded
path, \(H_j\) is exponential with mean \(1/\Lambda(X_j)\).  Tonelli's theorem
and conditional expectation give
\[
 \mathbb E\sum_{j=0}^{N-1}H_j
 =\mathbb E\sum_{j\ge0}\mathbf1_{\{j<N\}}
   \frac1{\Lambda(X_j)}
 \le\frac{\mathbb EN}{\inf_x\Lambda(x)}<\infty.
\]
Thus the expected physical return time is finite, provided the construction
is nonexplosive.

Finite mean return makes \(x_*\) recurrent for the embedded chain, so it is
visited infinitely often almost surely.  On each visit, the following
holding time has the same exponential law of finite rate \(\Lambda(x_*)\).
The strong Markov construction makes these recurrent holding times
independent copies.  Their infinite sum diverges almost surely.  Since total
physical time dominates this subseries, jump times cannot accumulate at a
finite time.  The minimal CTMC is nonexplosive and has finite mean positive
return time to \(x_*\).
\end{proof}

For the marked reaction chain, the carried target gives the required lower
rate bound at every augmented state:
\begin{equation}\label{eq:rate-lower}
  \Lambda(x)\ge\bar\kappa_t(x)_t\ge\bar\kappa_->0.
\end{equation}
Indeed, \(t\) is enabled and its falling factorial is a positive integer.
Propositions~\ref{prop:trace} and \ref{prop:ctmc} therefore prove
Theorem~\ref{thm:main}.  The same return-cycle construction yields the
stationary distribution in Corollary~\ref{cor:stationary}; uniqueness follows
from irreducibility.  Explicitly, after choosing a recurrent state \(x_*\),
normalizing the expected occupation measure during one return cycle gives a
stationary probability on \(\Gamma\).

\section{Scope and further questions}\label{sec:scope}

The proof is class-wise.  Coordinate faces, siphons, parity restrictions, and
other lattice constraints require no separate interior argument: every
constructed target-following path consists of actual enabled channels and
therefore remains in the fixed closed class.  Species absent from all
complexes are constant on each class and appear automatically in the
invariant branch of Lemma~\ref{lem:top}.

The one-linkage hypothesis is used at one decisive point: for every carried
target \(t\) and selected terminal complex \(c\), a directed path
\eqref{eq:path} must exist.  With several linkage classes, a carried target
need not reach a terminal source belonging to another class, so the finite
path library does not directly glue.  That extension remains open here.

Bimolecularity is used in the exhaustive case split of
Lemma~\ref{lem:top}: a top complex either contains two diverging particles or
exactly one, leaving at most one service species.  Complexes of molecularity
three or higher admit additional top configurations not covered by the
argument.  No claim is made for them.

The proof establishes positive recurrence and the existence of a stationary
probability distribution.  It does not provide a product form, explicit tail
bounds, exponential ergodicity, or quantitative mixing rates.

\section*{Declarations}

\paragraph{Author information.}
Alec Kriebel, Independent researcher.  Version 0.2, 6 August 2026.\\
Correspondence: \mbox{\href{mailto:me@aleckriebel.com}{me@aleckriebel.com}}.

\paragraph{Code availability.}
A self-contained verification archive accompanies this manuscript.  It
contains deterministic channel-level identity checks, adversarial examples,
and fixed-seed stress tests.  No computation is load bearing for the
universal proof.

\paragraph{Funding.}
This research received no specific grant from any funding agency in the public,
commercial, or not-for-profit sectors.

\paragraph{Competing interests.}
The author declares no competing interests.

\subsection*{A note on the use of AI}
OpenAI ChatGPT, model GPT-5.6 Pro, accessed through the ChatGPT interface on
5--6 August 2026, was used extensively for exploratory proof search,
generation and testing of verification code, drafting, and adversarial
editorial critique.  The interaction included attempts to construct
counterexamples and several proposed approaches that were later rejected,
including universal one-step Lyapunov claims and rate-changing cycle
comparisons.  The final manuscript uses only the marked-channel factorial
argument reconstructed in Sections~\ref{sec:marked}--\ref{sec:recurrence}.

The author reviewed the released manuscript and verification outputs and
assumes responsibility for all mathematical claims, citations, code, and
conclusions.  Automated checks are not peer review, and this declaration does
not assert that every claim has already received independent verification by
another human expert.  The AI system is not listed as an author.

\bibliographystyle{plainnat}
\bibliography{references}

\appendix
\section{Audit details for the trace chain and physical time}\label{app:trace}

This appendix records several interfaces that can otherwise be obscured by
standard Markov-chain shorthand.

First, the augmented state records the actual marked channel.  If two
channels have the same population displacement but different sources or
targets, their post-jump populations can agree while their carried targets do
not.  They therefore remain separate.  Only channels with identical source
and target may be aggregated.

Second, \(K\) in \eqref{eq:K} is nonempty, not merely finite.  The global
minimum argument in Proposition~\ref{prop:K} is pathwise: every possible
episode endpoint has potential at least the starting minimum.  No optional
stopping or integrability assertion is needed for this step because each
episode has at most \(|\mathcal C|\) jumps.

Third, the trace-chain proof distinguishes hitting from return.  The first
ordinary jump from \(k\in K\) forces a positive return time.  The finite
expected time back to \(K\) follows by conditioning on finitely many
successors.  The subsequent geometric-block argument is applied to the
finite trace chain, and only afterward are the random trace excursions
expanded into original embedded jumps.

Fourth, nonexplosion is not inferred from recurrence without an argument.
The embedded chain is constructed first.  Its recurrent marked state is
visited infinitely often.  The holding times attached to those visits form
an infinite sequence of independent exponentials with one fixed finite rate,
whose sum diverges almost surely.  This directly prevents finite-time
accumulation.

Finally, singleton absorbing classes are handled before augmentation.  In an
infinite irreducible class every population state has an enabled genuine
channel; otherwise it would be absorbing and its class would be a singleton.

\section{Computational verification}\label{app:verification}

The accompanying Python package has no runtime third-party dependencies.  It
checks the exact falling-factorial identity using rational arithmetic, checks
the entropy rewrite by representing logarithms of positive rationals as
rational prime-exponent signatures, verifies the scalar-envelope branch
conditions exactly, and tests the top-complex classification on a
deterministic finite atlas.  It also includes fixed-seed random stress tests
and the following adversarial calibrations:

\begin{itemize}[leftmargin=2em,itemsep=0pt,topsep=0.1em,parsep=0pt,partopsep=0pt]
\item the zero complex as source and target;
\item pure binary and mixed complexes \(2A\) and \(A+B\);
\item permanently zero coordinates on a closed face;
\item a diverging species with normalized logarithmic weight zero;
\item species absent from every complex;
\item shared service species;
\item parallel channels and distinct channels with one displacement;
\item the zero-length path \(c=t\);
\item equal-molecularity and absorbing finite classes;
\item signed linear invariants and parity-restricted paths;
\item the cycle \(0\to A+B\to B\to0\).
\end{itemize}

Finite enumeration and random testing are adversarial calibration only; they
are not substitutes for Lemma~\ref{lem:top}.  The canonical verification JSON
excludes elapsed times, timestamps, temporary paths, and platform-dependent
formatting, and is required to agree byte for byte across two runs.
